\section{The Hunting of the Symplectic Phantasm}
\label{sec:symplectic-phantasm}

The direction adopted on 9 September 2026 begins with finite symplectic
objects and their quantum realizations, and then asks how arithmetic
processes can act on systems assembled across primes. The procedure is
sober, accretive, and careful: primary sources precede definitions,
definitions precede lemma contracts, and proofs precede promotion.
The relation to the Riemann hypothesis remains a motivation.

The following definitions fix a first layer. All propositions in this
section are sketches of obligations to be discharged, with no new claim
promoted to proved status. In particular, a category of quantum processes,
a concrete observable algebra, and a representation with a reference state
are separate objects. Their comparison is part of the work.

\subsection{Definition layer}

\begin{definition}[finite symplectic objects and affine symmetries]
Fix a finite field $k$ of characteristic $p$. A symplectic object is a
finite-dimensional $k$-vector space $V$ with a $k$-bilinear alternating form
$\omega_V$ whose radical is zero. Its rank is $n$ when $\dim_k V=2n$.
The zero space is allowed. Write $\overline V=(V,-\omega_V)$ and equip
$V\oplus W$ with $\omega_V\oplus\omega_W$. A subspace $L\subseteq V$ is
Lagrangian when $L=L^\perp$, where
$L^\perp=\{v:\omega_V(v,L)=0\}$.
The linear symmetry groupoid $\mathsf S_k$ has these objects and
form-preserving linear isomorphisms. The affine symmetry groupoid
$\mathsf S_k^{\mathrm{aff}}$ has arrows $(t,g):V\to W$, where
$g:V\to W$ is such an isomorphism and $t\in W$, acting by $v\mapsto gv+t$.
Composition is $(s,h)\circ(t,g)=(s+ht,hg)$ and the identity is $(0,1_V)$.
\end{definition}
\begin{scope}
All characteristics are included in this classical datum. Rank one means dimension two; no quantum realization is part of the definition.
\end{scope}
\provenance{D1701}{Stipulation; SP-GH07, SP-W09}{See the lemma contracts below}{Source-first bootstrap; no promotion}

\begin{definition}[affine Lagrangian relations]
Fix a finite field $k$ and finite-dimensional symplectic $k$-spaces $V,W$.
An arrow $R:V\to W$ is either the empty relation or an affine translate
of a Lagrangian linear subspace of $\overline V\oplus W$.
For $R:V\to W$ and $S:W\to Z$ prescribe
\[
 S\circ R=\{(v,z):\text{there exists }w\in W\text{ with }(v,w)\in R,
 (w,z)\in S\}.
\]
The identity is $\Delta_V=\{(v,v):v\in V\}$; the dagger is relational
converse. The product is the direct sum on objects and the Cartesian
product on relations, with coordinates reordered from
$(\overline V\oplus W)\oplus(\overline {V'}\oplus W')$ to
$\overline{V\oplus V'}\oplus(W\oplus W')$. The unit object is the zero
symplectic space. Write $\mathsf L_k^{\mathrm{aff}}$ for this candidate.
\end{definition}
\begin{scope}
Closure and category/coherence laws are proof obligations. Empty relations are retained; nonlinear subvarieties are outside this definition.
\end{scope}
\provenance{D1702}{Stipulation; SP-W09, SP-LW14, SP-CK21}{See the lemma contracts below}{Source-first bootstrap; no promotion}

\begin{definition}[odd-characteristic Weyl datum in arbitrary rank]
Fix a finite field $k$ of odd characteristic, a nontrivial additive
character $\psi:k\to U(1)$, and a finite-dimensional symplectic
$k$-space $(V,\omega_V)$. Put $\beta_V=\omega_V/2$.
The formal Weyl algebra $\mathcal A_{k,\psi}(V)$ has complex basis
$\{w_v:v\in V\}$ and prescribed operations
\[
 w_vw_z=\psi(\omega_V(v,z)/2)w_{v+z},\qquad
 w_v^*=w_{-v},\qquad 1=w_0.
\]
Prescribe $\tau_V(\sum_v c_vw_v)=c_0$. On $k\times V$ prescribe the
Heisenberg multiplication
$(a,v)(b,z)=(a+b+\omega_V(v,z)/2,v+z)$, with central character $\psi$.
For the standard space $V=k^n\oplus k^n$ with
$\omega_V((a,b),(a',b'))=a\cdot b'-a'\cdot b$, put
$\mathcal H_{k,n}=\ell^2(k^n)$ with counting inner product and prescribe
\[
 (W(a,b)f)(x)=\psi(b\cdot x+a\cdot b/2)f(x+a).
\]
An abstract-space model additionally names a symplectic identification
with the standard space. Model independence is a claim, not a hidden choice.
\end{definition}
\begin{scope}
The half-form convention excludes characteristic two. Algebra laws, positivity, matrix realization and uniqueness are separate obligations.
\end{scope}
\provenance{D1703}{Stipulation; SP-GH07, SP-PRASAD09, SP-GROSS06}{See the lemma contracts below}{Source-first bootstrap; no promotion}

\begin{definition}[actual stabilizer amplitudes]
Fix an odd prime $p$, the character $\psi_p(x)=\exp(2\pi i x/p)$, and
$\mathcal H_{p,n}=\ell^2(\mathbb F_p^n)$ for $n\geq0$.
Let $\mathcal P_{p,n}$ be the subgroup of unitaries generated by the
standard Weyl operators and all scalar phases. A Clifford unitary is a
unitary $U$ satisfying $U\mathcal P_{p,n}U^*=\mathcal P_{p,n}$.
Define $\mathsf{Stab}^{\mathrm{amp}}_p$ to have objects
$\mathcal H_{p,n}$ and actual linear maps generated under composition,
tensor product and adjoint by Clifford unitaries, the map
$\eta:\mathbb C\to\mathcal H_{p,1}$ with $\eta(1)=\delta_0$,
and scalar maps $c:\mathbb C\to\mathbb C$ for every $c\in\mathbb C$.
Use the standard coordinate identification
$\mathcal H_{p,m+n}=\mathcal H_{p,m}\otimes\mathcal H_{p,n}$.
Equality is equality of the resulting linear maps, including zero.
\end{definition}
\begin{scope}
Scalars and amplitudes are retained. Such a map defines a physical successful branch only with the required norm or instrument normalization. No closure under addition is imposed here.
\end{scope}
\provenance{D1704}{Stipulation; SP-GROSS06, SP-CK21}{See the lemma contracts below}{Source-first bootstrap; no promotion}

\begin{definition}[stabilizer amplitudes modulo invertible scalars]
Fix an odd prime $p$ and the actual stabilizer amplitude category with
objects $\ell^2(\mathbb F_p^n)$. In each Hom-set identify two nonzero
maps $T,S$ precisely when $S=cT$ for some $c\in\mathbb C^\times$.
Keep the zero map in a separate class. Write
$\mathsf{Stab}^{\mathrm{proj}}_p$ for the quotient with the induced
composition, tensor product and adjoint. This scalar quotient is different
from quotienting unitary intertwiners only by $U(1)$.
\end{definition}
\begin{scope}
The quotient records no success probability or norm. A coherent sum of independently rescaled representatives is not prescribed by this definition.
\end{scope}
\provenance{D1705}{Stipulation; SP-CK21}{See the lemma contracts below}{Source-first bootstrap; no promotion}

\begin{definition}[finite quantum branches and instruments]
A finite quantum system is a finite nonempty family of nonzero
finite-dimensional complex Hilbert spaces $\mathbf H=(H_a)_{a\in A}$
with algebra $\mathcal B(\mathbf H)=\bigoplus_a\operatorname{End}(H_a)$
and trace $\operatorname{Tr}_{\mathbf H}=\sum_a\operatorname{Tr}_{H_a}$.
A branch $\Phi:\mathcal B(\mathbf H)\to\mathcal B(\mathbf K)$ is a
linear map admitting finitely many Kraus operators $K_{bar}:H_a\to K_b$
with
\[
 \Phi(x)_b=\sum_{a,r}K_{bar}x_aK_{bar}^*,\qquad
 \sum_{b,r}K_{bar}^*K_{bar}\leq1_{H_a}\quad\text{for every }a.
\]
Branches are identified by equality of their maps. A channel has equality
in every displayed inequality. An instrument is a finite family of
branches $(\Phi_o)_o$ whose sum is a channel. Outcomes remain named.
For a positive input $\rho$ of trace one, define the probability of $o$
as $\operatorname{Tr}(\Phi_o(\rho))$ and its conditional state by division
by that probability only when it is positive. Use composition of maps,
tensor product of systems and maps, and the identity map.
\end{definition}
\begin{scope}
This is an ambient finite quantum process definition. It does not assert that every branch comes from the arithmetic or stabilizer generators. Traces are ordinary, not dimension-normalized.
\end{scope}
\provenance{D1706}{Stipulation; SP-WAT18}{See the lemma contracts below}{Source-first bootstrap; no promotion}

\begin{definition}[coherent additive completion and classical tags]
Fix an odd prime $p$ and its actual stabilizer amplitude category.
First take the complex linear span of each Hom-set inside the ambient
linear maps. Its matrix completion has objects finite lists
$(\mathcal H_{p,n_1},\ldots,\mathcal H_{p,n_s})$, including the empty
list, and arrows matrices whose entries lie in these linear spans.
Composition is matrix multiplication and dagger is adjoint transpose.
Realize a list as $H=\bigoplus_{i=1}^s\mathcal H_{p,n_i}$.
Its coherent observable algebra is $\operatorname{End}(H)$.
For a nonempty list the tagged observable algebra is instead
$\bigoplus_i\operatorname{End}(\mathcal H_{p,n_i})$, embedded as the
block-diagonal operators on $H$. Both constructions and their chosen
block decomposition are recorded; they are not identified with one another.
\end{definition}
\begin{scope}
This completion deliberately adds linear combinations. Whether it preserves a proposed Gaussian fragment, or comes from classical disjoint unions, is a separate question.
\end{scope}
\provenance{D1707}{Stipulation; SP-CK21, SP-WAT18}{See the lemma contracts below}{Source-first bootstrap; no promotion}

\begin{definition}[bosonic Fock space and bounded second quantization]
For a complex Hilbert space $H$, let $P_r=(r!)^{-1}\sum_{\pi\in S_r}U_\pi$
be the symmetrizer on the Hilbert tensor power $H^{\otimes r}$ and set
$\operatorname{Sym}^rH=\operatorname{ran}P_r$, with
$\operatorname{Sym}^0H=\mathbb C$. Define
\[
 \Gamma_s(H)=\widehat{\bigoplus}_{r\geq0}\operatorname{Sym}^rH,
 \qquad\Omega=1\in\operatorname{Sym}^0H.
\]
The hat means Hilbert direct sum. The algebraic finite-particle space is
$\bigoplus_{r\geq0}^{\mathrm{alg}}\operatorname{Sym}^rH$.
For a contraction $T:H\to K$, prescribe
$\Gamma_s(T)=\bigoplus_{r\geq0}T^{\otimes r}|_{\operatorname{Sym}^rH}$.
Particle number is the operator $N_H\xi=(r\xi_r)_r$ on the domain
$\{\xi:\sum_r r^2\|\xi_r\|^2<\infty\}$.
\end{definition}
\begin{scope}
Boundedness and functor laws are obligations. No bounded second quantization of an arbitrary expansive map is asserted, and no unspecified categorical free-monoid property is imposed.
\end{scope}
\provenance{D1708}{Stipulation; SP-DER06}{See the lemma contracts below}{Source-first bootstrap; no promotion}

\begin{definition}[scalar restriction and a named Frobenius datum]
Let $E/K$ be a finite extension of finite fields and let $(V,\omega_E)$
be a finite-dimensional symplectic $E$-space. On the underlying $K$-space
$\operatorname{Res}_{E/K}V$ prescribe
$\omega_K(v,w)=\operatorname{Tr}_{E/K}(\omega_E(v,w))$.
For a nontrivial additive character $\psi_K$, prescribe
$\psi_E=\psi_K\circ\operatorname{Tr}_{E/K}$.
A Frobenius datum consists of the automorphism
$\sigma(x)=x^{|K|}$ of $E$ and a bijective $\sigma$-semilinear map
$F:V\to V$ satisfying
$\omega_E(Fv,Fw)=\sigma(\omega_E(v,w))$.
On the standard space $E^n\oplus E^n$ choose coordinatewise $\sigma$.
On $\ell^2(E^n)$ its specified permutation is
$U_F\delta_x=\delta_{\sigma(x)}$.
\end{definition}
\begin{scope}
The classical data allow characteristic two; the half-form quantum comparison is restricted to odd characteristic. Abstract symplectic spaces do not acquire descent data by omission. This does not define a field-inclusion functor.
\end{scope}
\provenance{D1709}{Stipulation; SP-STFIELD, SP-GH07}{See the lemma contracts below}{Source-first bootstrap; no promotion}

\begin{definition}[a specified quantum subsystem decoder]
Fix an odd-characteristic finite field $k$, a nontrivial additive
character, and a symplectic linear injection $j:U\to V$ between
finite-dimensional symplectic $k$-spaces. Put $W=j(U)^\perp$.
A quantum subsystem datum additionally specifies finite-dimensional
Weyl model Hilbert spaces $H_U,H_W,H_V$ and a unitary
$J:H_U\otimes H_W\to H_V$ satisfying
\[
 J(W_U(u)\otimes W_W(w))J^*=W_V(ju+w).
\]
Define the observable inclusion $\iota_J(a)=J(a\otimes1)J^*$ and the
state decoder $\mathcal D_J(\rho)=\operatorname{Tr}_{H_W}(J^*\rho J)$,
where the partial trace is specified by
$\operatorname{Tr}(\mathcal D_J(\rho)a)=\operatorname{Tr}(\rho\iota_J(a))$
for every $a\in\operatorname{End}(H_U)$.
\end{definition}
\begin{scope}
The existence of the decomposition and compatible model unitary is an obligation. The datum is not scalar restriction, a generic field trace, or a non-bijective interpretation of finite-field Frobenius.
\end{scope}
\provenance{D1710}{Stipulation; SP-STFIELD, SP-WAT18}{See the lemma contracts below}{Source-first bootstrap; no promotion}

\begin{definition}[finite-prime tensor assembly with a chosen reference]
For every prime $p$, choose an integer $d_p\geq1$, a Hilbert space
$H_p=\mathbb C^{d_p}$ and a positive operator $\rho_p$ of trace one.
For a finite set of primes $P$, put
$A_P=\bigotimes_{p\in P}\operatorname{End}(H_p)$, using increasing prime
order, and $A_\varnothing=\mathbb C$. For $P\subseteq Q$ prescribe
$\iota_{QP}(a)=a\otimes1$ with the required coordinate reordering.
Define the candidate norm completion
$A_{\mathrm{pr}}=\overline{\varinjlim_P A_P}^{\|\cdot\|}$ using the
finite matrix norms. On finite tensors prescribe
$\varphi_{\mathrm{pr}}(\bigotimes_{p\in P}a_p)
=\prod_{p\in P}\operatorname{Tr}(\rho_pa_p)$.
\end{definition}
\begin{scope}
The inductive and state-extension properties belong to a lemma. The matrices are specified data, not yet an all-prime Weyl assignment. Neither inter-prime arithmetic maps nor a factor type is supplied.
\end{scope}
\provenance{D1711}{Stipulation; SP-CM08, SP-WAT18}{See the lemma contracts below}{Source-first bootstrap; no promotion}

\begin{definition}[a represented Bost--Connes control system]
On $H_{\mathrm{BC}}=\ell^2(\mathbb N_{>0})$ with basis $(\delta_m)_{m\geq1}$,
fix the usual embedding $\mathbb Q/\mathbb Z\to U(1)$ given by
$r\mapsto\exp(2\pi i r)$, and prescribe
\[
 e(r)\delta_m=\exp(2\pi i mr)\delta_m,\qquad
 \mu_n\delta_m=\delta_{nm}\quad(n\geq1).
\]
Let $A_{\mathrm{BC}}^{\mathrm{rep}}=C^*(e(r),\mu_n:r\in\mathbb Q/\mathbb Z,
n\geq1)\subseteq B(H_{\mathrm{BC}})$.
Define $H_{\log}\delta_m=(\log m)\delta_m$ on
$\{\xi:\sum_m(\log m)^2|\xi_m|^2<\infty\}$ and prescribe
$\sigma_u(a)=e^{iuH_{\log}}ae^{-iuH_{\log}}$.
Define $\nu_n(a)=\mu_na\mu_n^*$ and $L_n(a)=\mu_n^*a\mu_n$.
For real $b>1$ prescribe
$Z_{\mathrm{BC}}(b)=\sum_{m\geq1}m^{-b}$ and
$\rho_{\mathrm{BC},b}=Z_{\mathrm{BC}}(b)^{-1}e^{-bH_{\log}}$.
\end{definition}
\begin{scope}
This is a concrete represented control, not an assertion of faithfulness of a universal presentation or an identification with the sought symplectic system. Convergence, invariance and CP properties remain lemma obligations.
\end{scope}
\provenance{D1712}{Stipulation; SP-BC95, SP-CM04, SP-CM08}{See the lemma contracts below}{Source-first bootstrap; no promotion}

\begin{definition}[state, GNS and modular comparison data]
A $C^*$-dynamical datum is a unital $C^*$-algebra $A$ with a point-norm
continuous homomorphism $\sigma:\mathbb R\to\operatorname{Aut}(A)$.
A state is a positive linear functional $\varphi$ with $\varphi(1)=1$.
Its GNS prescription is the completion of $A/N_\varphi$ with
$N_\varphi=\{a:\varphi(a^*a)=0\}$ and
$\langle[a],[b]\rangle=\varphi(a^*b)$, left action
$\pi_\varphi(a)[b]=[ab]$, and cyclic vector $\Omega_\varphi=[1]$.
Write $M_\varphi=\pi_\varphi(A)''$.
For real $b>0$, a $\mathrm{KMS}_b$ state means that for every $a,c\in A$
there is a bounded continuous function on $0\leq\operatorname{Im}z\leq b$,
holomorphic in the interior, whose boundary values are
$F_{a,c}(u)=\varphi(a\sigma_u(c))$ and
$F_{a,c}(u+ib)=\varphi(\sigma_u(c)a)$.
A modular comparison datum additionally requires $\Omega_\varphi$ to be
cyclic and separating for $M_\varphi$. With
$S_\varphi(a\Omega_\varphi)=a^*\Omega_\varphi$ and polar decomposition
$\overline S_\varphi=J_\varphi\Delta_\varphi^{1/2}$, its modular-flow
prescription is $a\mapsto\Delta_\varphi^{iu}a\Delta_\varphi^{-iu}$.
\end{definition}
\begin{scope}
Analytic existence and comparison theorems must be cited or proved before use. No equality between modular flow, the specified physical flow, or Frobenius is a stipulation.
\end{scope}
\provenance{D1713}{Stipulation; SP-CM08, SP-CCM07}{See the lemma contracts below}{Source-first bootstrap; no promotion}
